%── PUBLICATIONS & WORKING PAPERS ────────────────────────────────────────────
\section{Publications \& Working Papers}

\cvpub{\href{https://www.sciencedirect.com/science/article/abs/pii/S0957178725000189}{Private Ownership of Water and Wastewater Systems: Assessing Health Impacts}}
      {Utilities Policy, peer-reviewed journal article \textnormal{(first author)}}
      {In 2020 Brazil made it easier for private companies to run water and sewage systems. This paper asks whether that was a good idea by looking at the municipalities that had already privatised. Comparing closely matched municipalities before and after they switched, over 1998 to 2021, we find fewer hospital admissions for diarrhoeal disease among children under five (ATT $-11.6$ for infants, $-15.8$ at ages 1 to 5) and no change in diseases unrelated to sanitation. Results vary across municipalities, and the clearest evidence comes from Tocantins, the one state that moved to full private ownership. I wrote the paper, built the dataset, and wrote the code; my co-author supervised.}

\cvpub{\href{https://papers.ssrn.com/sol3/papers.cfm?abstract_id=5441274}{Resistance to Doing Right: Rising Cost of Misconduct from Body-Worn Camera Adoption}}
      {Working paper, in progress \textnormal{(second author)}}
      {Between 2021 and 2022 São Paulo issued 10,125 body cameras across 64 police battalions. We use confidential Military Police records from 2020 to 2023 to ask what changed. Battalion boundaries are not public, so I rebuilt them by clustering the locations of recorded incidents into Voronoi polygons, which is what made it possible to say which areas were treated. Confrontations between police and civilians fell (ATT $-0.50$, significant at 1\%), with the drop concentrated in low-income areas (5\%, and no significant change in high-income ones), while property crime rose, most clearly vehicle theft. That pattern looks less like better policing and more like less policing.}
